\textbf{T2: Hysteresis (10 pts)}

Consider a solenoid with $N \gg 1$ turns, radius $R$ and length
$\ell \gg R$. The solenoid has a ferromagnetic core with hysteresis curve as
shown in the picture. $M$ denotes the magnetization of the core and would be
zero without core material. It is related to $H$, the magnetic field strength
generated by the current in the solenoid coil, and $B$ by
$B = \mu_0\,(H + M)$. $M_0$ and $H_0$ are of the same order of magnitude.
Assume that $M$ can only change if $|H| \geq H_0$.

\begin{center}\includegraphics[width=0.42\linewidth]{figures/EuPhO_2026_Q2_fig1.png}\end{center}

An ideal capacitor with capacitance $C$ is now connected to the solenoid
forming a closed circuit. Assume that all wires have negligible resistance.

a) \textbf{(3.0 pts)} Initially there is a current $I_\mathrm{i}$ flowing
through the circuit and the capacitor is uncharged. $I_\mathrm{i}$ is large
enough that $H(I_\mathrm{i}) \gg H_0$. After one oscillation the current again
reaches a maximum value. Determine the difference between $I_\mathrm{i}$ and
this value.

b) \textbf{(2.3 pts)} Find the maximum current that can be attained after many
oscillations.

c) \textbf{(4.7 pts)} The current $I(t)$ exhibits two qualitatively distinct
phases of behaviour, A and B. The system can remain in phase A indefinitely,
whereas any single interval spent in phase B has bounded duration. Find the
maximal possible duration of a phase-B interval, if $I_\mathrm{i}$ is chosen
optimally.
